\chapter{Attitude Simulation}

To demonstrate the performance of the FNDO-BS controller derived in chapter \ref{chap:underactuated_controller}, it is necessary to simulate it's performance in a realistic environment. In this chapter, building on the work of the previous chapters,the performance of the FNDO-BS controller is compared to a baseline PID controller during a realistic operation scenario. In previous chapters, the EOSSAT-1 satellite structure was tested against.  The performance of the controllers is evaluated based on their ability to maintain the desired attitude and minimize the control effort required to do so. A realistic simulation environment is created using the parameters derived in chapter \ref{chap:modelling}, and the results of the simulations are presented in this chapter.

A simulation was conducted using the parameters derived in chapter \ref{chap:modelling}.

The simulation was conducted in python, and the results are presented in this chapter.

Experimental setup:

% @TODO: Explain tuning of PID controller for EOS.


\vspace{1cm}

\begingroup
\begin{table}
    \centering
    \begin{tabular}{@{}p{5cm} p{2cm} p{5cm} c}
        Parameter & Notation & Value & Unit \\
        \hline
        Moment of Inertia & $J_0$ & $\begin{bmatrix}
            2.5 & 0 & 0 \\
            0 & 3.5 & 0 \\
            0 & 0 & 4.5
        \end{bmatrix}$ & $kg\,m^2$  \\
        \hline
        Sampling time (controller) & $T_{s,ctrl}$ & 0.01 & s \\
        \hline
        Sampling time (MAGT) & $T_{s,mag}$ & 1.0 & s \\
        \hline
        
        \hline
    \end{tabular}
\caption{Table of reference frames shorthands used in this thesis.}
\end{table}
\endgroup

\vspace{1cm}

